\section{Implementation and calibration}
\label{app:details}

Store seeds are 0--11 for the count fit and 100--129 for the held-out test.
We sample three syllables from a fixed list of sixteen, rejecting duplicate keys and lines not matching the old rule's tokenized length.
All active and inactive entries have the same sentence shape except for the initial \texttt{Tickets} or \texttt{Notes} marker and the key.
The full prompt follows Qwen's chat markers and ends in an assistant cue, \texttt{Queue}, followed by a space.
Each candidate digit is one token following that cue.
We read six-way log probabilities with an equal logit offset on the six digits and retain exactly those six probabilities.
The offset and common normalizer cancel in pairwise log-odds.
The scorer checks the returned token identities and verifies invariance to a larger offset on the calibration prompt.

Query 1 asks where a newly arrived ticket belongs; query 2 asks which queue handles the ticket type. Both request only a number.
We use query 1 to check each new key in its fully active cube and score the shared all-inactive cell once per store and wording.
On the calibration prompt, prefix caching agrees with an uncached read within $10^{-3}$ log-probability units.
We run inference only on an eight-thread Intel Xeon CPU.

Bootstrap intervals use 10,000 resamples of the held-out store indices with replacement.
The exact one-sided zero-event upper limit, where applicable, solves $(1-p)^n=0.05$ for $n$ store draws.

\paragraph{Synthetic crossing checks.}
The additive margin $m(S)=2.5-|S|$ fails only at the full set, but its third difference is zero.
A count curve $(2.5,2,1,-1)$ also fails only at the full set and gives third difference $-0.5$ without any identity effect.
If all proper subsets instead score $1$ and the triple scores $-1$, the pair forecast is $1$ and misses by $-2$ logits (observed minus predicted).
These cubes check the analysis, not the language model.
